\begin{lemma}\label{lem:decic60-quintic-descent}
Let $F$ be a quintic number field, and suppose that $E/F$ has a point
of order $30$. Put $L=F(E[2])$, and suppose that $[L:F]=2$.
No curve over $L$ obtained from $E$ by a $2$-isogeny has a point
of order $60$.
\end{lemma}
\begin{proof}
The classification of quintic points on $X_1(30)$ gives the two
curves $E_{30},E_{30}'$ over
$F=\mathbb Q(\alpha)$, where
$\alpha^5+\alpha^4-3\alpha^3+3\alpha+1=0$.
We use the printed Tate models in the quintic classification.
They are $2$-isogenous over $F$.
For each model, its $2$-division polynomial has one linear and
one irreducible quadratic factor $h$ over $F$.
The exact calculations in \path{quintic_2tower.m} give
\[
 -\operatorname{disc}(h)\notin F^{\times2}.
\]
Since $[F:\mathbb Q]$ is odd, $i\notin F$; hence $i\notin L$.
The same calculation shows that the polynomial for halving the
chosen point of order $30$ is irreducible of degree $4$ over $F$,
for both curves.

Both full $2$-division fields have a degree-one prime above $53$
at which the curve has good reduction with exactly $60$ points.
The explicit quintic and quadratic residue roots are recorded in
the pinned full-$2$ certificates. Thus torsion on every
$L$-rational $2$-isogenous curve has order dividing $60$.
In particular, $E(L)_{\mathrm{tors}}=C_2\times C_{30}$.
The quotient by the $F$-rational point of order $2$ is the other
quintic model. It has no point of order $4$ over $L$: otherwise its
$2$-primary torsion is cyclic of order $4$, and halving a generator
of its $F$-rational subgroup of order $30$ would be possible over
$L$, contrary to the irreducible quartic.

Complete the square and write the roots of the $2$-division cubic
as $r_0,r_1,r_2$, with $r_0\in F$ and $r_1,r_2$ exchanged by the
nontrivial element $\sigma$ of $\operatorname{Gal}(L/F)$.
A point $(r_0,0)$ is divisible by $2$ precisely when both
$r_0-r_1$ and $r_0-r_2$ are squares. Since they are conjugate and
$E(L)$ has no point of order $4$, both are nonsquares.
The absence of order $4$ on the quotient by $(r_0,0)$ similarly
shows that $r_1-r_0$ and $r_2-r_0$ are nonsquares.
Here we use the usual $2$-isogeny formula: on the quotient by
$(r_j,0)$, the nonzero dual-kernel point is divisible by $2$ if and
only if at least one of the two differences $r_k-r_j$, $k\ne j$,
is a square.

Suppose that the quotient by $(r_1,0)$ has a point of order $4$.
Its $2$-primary torsion has order at most $4$, so that point halves
the dual-kernel point. The formula just recalled implies that
$r_2-r_1$ is a square, since $r_0-r_1$ is not.
Applying $\sigma$ shows that its negative is also a square.
This gives $i\in L$, a contradiction. The quotient by $(r_2,0)$
is conjugate and is excluded as well. None of the three quotients
has a point of order $4$, and hence none has a point of order $60$.
\end{proof}

\begin{proposition}\label{prop:decic-no60}
An elliptic curve over a number field of degree $10$ has no point
of order $60$.
\end{proposition}
\begin{proof}
Suppose that such a point exists. Its field of moduli has degree
$10$, since orders $60$ and $C_2\times C_{30}$ are excluded in
all proper divisor degrees $1,2,5$ by the known classifications.
Quotienting by the rational subgroup of order $2$ gives a point
on $X=X_1(2,30)$ over the same field $K$.
Put $E_0=X_1(15)$ and let $Y$ be the quotient of $X$ by the
alternating subgroup on the three nonzero $2$-torsion points.
Then
\[
 X\longrightarrow Y\xrightarrow{f}E_0,
 \qquad [\mathbb Q(Y):\mathbb Q(E_0)]=2,
 \qquad \mathbb Q(Y)=\mathbb Q(E_0)(\sqrt{\Delta}).
\]
Here $\Delta$ is the discriminant of the universal Tate curve.
The first map is a $C_3$-cover away from the cusps, and $g(Y)=9$.
The resulting point on $Y$ also has degree $10$, since its residue
extension in the $C_3$-cover has degree $1$ or $3$.
Let $D_Y$ be its rational trace divisor.

The rational Jacobian of $X$ has rank zero and is generated by
rational cuspidal divisors
\cite[Theorem~3.1, Proposition~4.16 and Table~3]{DEHMZ}.
Consequently
$[D_Y-10c_0]$ belongs to the subgroup generated by the eight
rational cusp orbits $C_i$ on $Y$. Their degrees are
\[
 (1,1,4,2,4,1,2,1).
\]
They are the unique ramified lifts of the corresponding cusp
orbits on $E_0$. The modular-unit valuation signatures identify
these global orbits after reduction at $7$.
The proven finite presentation gives
\[
 J_Y(\mathbb F_7)=[2,2,2,2,2,2,102,10608],\qquad
 H=[2,2,2,2,2,8],\qquad \#H=256.
\]
Reduction identifies this cuspidal subgroup with its specialization.

We compute all $256$ normalized classes in degree $10$, each with
an effective cuspidal representative $L=\sum_i v_iC_i$.
The decomposition under the involution of $f$ is explicit.
If $\nu_i$ is the valuation of $\Delta$ at the corresponding base
cusp, its two eigenspaces are the section spaces on $E_0$ of
\[
 B_+=\sum_i\left\lfloor\frac{v_i}{2}\right\rfloor f(C_i),
 \qquad
 B_-=\sum_i\left\lfloor\frac{v_i+\nu_i}{2}\right\rfloor f(C_i).
\]
All $\nu_i$ are odd. Genus-one Riemann--Roch determines these
spaces over $\mathbb Q$: in degree zero, principality is preserved
by reduction because the divisors are cuspidal torsion.
The finite Riemann--Roch dimensions agree with the sum of these
two dimensions in every class. The exhaustive alternatives are
\[
\begin{array}{c|r}
\text{series with a positive fixed cuspidal part}&168\\
\text{free pullback series }f^*|B|,\ \deg B=5&4\\
\text{series with both eigenspaces nonzero}&84.
\end{array}
\]
In the first row, the pairs of eigenspace dimensions are
$(2,0)$, $(3,0)$, and $(4,0)$, occurring $72$, $64$, and $32$
times, respectively. Writing
$L=f^*B_++C$ with $C=\sum_i(v_i\bmod2)C_i$, the divisor $C$
has degree $6$, $4$, or $2$, respectively. Since the elliptic base
has exactly $\deg B_+$ sections in positive degree, the
characteristic-zero lower bound equals the finite upper bound.
Hence $|L|=f^*|B_+|+C$ over $\mathbb Q$, and $C$ is a positive
fixed global cusp divisor. The first row cannot contain $D_Y$.

Suppose that $D_Y$ lies in the second row.
Since $D_Y$ is an irreducible divisor of degree $10$, its image on
$E_0$ is a point of degree $5$. The fine $15$-level moduli problem
therefore gives a model $E'/F$ with a point of order $15$, where
$[F:\mathbb Q]=5$ and $[K:F]=2$.
The full $2$-torsion of $E'$ is defined over $K$.
Its $2$-division cubic thus has an $F$-rational root, since an
irreducible cubic cannot split in a quadratic extension.
The quintic classification excludes full $2$-torsion with a point
of order $30$ over $F$, so $K=F(E'[2])$.
Adding the rational point of order $2$ to the point of order $15$
gives an $F$-rational point of order $30$.
The original curve is one of the three $2$-isogenous quotients of
$E'_K$, contrary to Lemma~\ref{lem:decic60-quintic-descent}.

It remains to treat the $84$ mixed series.
They consist of $60$ pencils and $24$ spaces of dimension $3$.
We enumerate all their $1848$ projective members over $\mathbb F_7$
and every oriented full-$2$ lift to $X$.
An irreducible $2$-division cubic requires a support multiplicity
divisible by $3$; when the cubic splits, every distribution among
the three oriented lifts is included.
The resulting $3414$ possible noncuspidal reductions involve
$1674$ distinct source places. Every cuspidal complement is left
unrestricted.

Apply the three maps to $X_1(30)$ retaining
$Q$, $P_2+16Q$, and $P_2+Q$.
Its rational Jacobian is cuspidal with invariants $[4,8160]$, and
its rational image in the proven finite Picard group has index
$4352$. We allow every finite target cusp complement of the
required degree, with all residue norms.
Exact addition of the three normalized classes reduces the
$3414$ patterns to $24$: twelve of type $2+2+3+3$, and twelve of
type $2+4+4$. All have noncuspidal degree $10$.

For each of these $24$ patterns, at least one multiplicity-one
support has no lift to the original order-$60$ problem.
To check this, let $(E',P_2,Q)$ represent that support, form
$\phi:E'\to E'/\langle P_2\rangle$, and let $\widehat\phi$ be
its dual. We enumerate the halves of $\phi(Q)$ and retain precisely
those $P$ satisfying $\widehat\phi(P)=Q$.
There are zero or two such points over the original residue field.
A multiplicity-one support cannot use an inert quadratic lift.
The exact tests on all $72$ active source places remove every
remaining pattern.

Finally, entirely cuspidal reduction of the original trace on
$X_1(60)$ at $7$ is impossible: rank zero and Manin--Drinfeld make
it linearly equivalent over $\mathbb Q_7$ to its cuspidal lift,
whereas
$\operatorname{gon}_{\mathbb Q}X_1(60)\geq7\cdot1152/800>10$.
Base change of $H^0$ gives the contradiction even for a local
cuspidal lift. Thus all reductions have been covered.
\end{proof}
